\documentclass[11pt]{article}
\usepackage[margin=1.0in]{geometry}
\usepackage{amsmath,amssymb}
\usepackage{booktabs}
\usepackage{graphicx}
\usepackage{caption}
\usepackage{longtable}
\usepackage[colorlinks=true,linkcolor=blue,citecolor=blue,urlcolor=blue]{hyperref}
\usepackage{titlesec}
\usepackage{enumitem}
\usepackage{fancyhdr}
\usepackage{xcolor}
\usepackage{mdframed}
\usepackage{url}
\urlstyle{tt}
% Long repo paths must be allowed to break, or they run off the page.
\newcommand{\fp}[1]{\path{#1}}
\sloppy

\captionsetup{font=small,labelfont=bf}
\titleformat{\section}{\normalfont\large\bfseries}{\thesection}{0.6em}{}
\titleformat{\subsection}{\normalfont\normalsize\bfseries}{\thesubsection}{0.6em}{}
\titlespacing*{\section}{0pt}{10pt}{5pt}
\setlist{itemsep=2pt,parsep=0pt,topsep=3pt}

\definecolor{boxbg}{RGB}{245,245,245}
\newmdenv[backgroundcolor=boxbg,linewidth=0.4pt,linecolor=gray,
          innertopmargin=6pt,innerbottommargin=6pt,skipabove=6pt,skipbelow=6pt]{keybox}

\pagestyle{fancy}
\fancyhf{}
\lhead{\small Kahlus --- primer}
\rhead{\small 28 July 2026}
\cfoot{\thepage}
\renewcommand{\headrulewidth}{0.4pt}

\begin{document}

\begin{center}
{\LARGE\bfseries Kahlus: a primer}\\[6pt]
{\large Everything the project rests on, explained from zero}\\[12pt]
{\small Written to be read straight through. No prior background assumed.}
\end{center}

\vspace{4pt}\hrule\vspace{8pt}

\begin{keybox}
\textbf{How to use this.} Part I is what we are actually researching and why.
Part II is the vocabulary --- if a word in a meeting confuses you, it is defined
there. Part III is the statistics. Part IV explains what ``F9'' means and lists
every finding. Part V works through the two real results in full, equation by
equation. Part VI is every paper we rely on and what each one is for. Part VII is
a glossary, Part VIII is the questions you will get asked.

\textbf{Where I am uncertain I say so.} Anything marked \textbf{[ours]} is a claim
this project makes; anything marked \textbf{[established]} is standard published
background you can assert freely.
\end{keybox}

\section{Part I --- What are we actually researching?}

\subsection{The question in one sentence}

\textbf{Can you see a change in someone's brain state before it happens?}

Not diagnose it after. Not detect it as it occurs. \emph{Predict} it, from the
signal alone, with enough lead time to do something about it. That is the
long-term goal --- seizures, arousals from sleep, lapses in attention.

Everything in this project is machinery for answering a narrower, more honest
version of that question: \textbf{does a brain signal contain information about
its own future that a simple method cannot already extract?} If a complicated
model does no better than a simple one, the complicated model has found nothing.

\subsection{What EEG is} \textbf{[established]}

Electroencephalography. Electrodes on the scalp measure voltage fluctuations
caused by the synchronised electrical activity of large populations of cortical
neurons --- principally the summed postsynaptic potentials of pyramidal cells
aligned perpendicular to the cortical surface. The signal is tiny, on the order of
tens of microvolts, and it is filtered by the skull and scalp, so it is spatially
blurred: one electrode sees a smeared average of a large patch of cortex.

Two properties matter enormously for us:

\begin{itemize}
\item \textbf{EEG is heavily autocorrelated.} The voltage 10\,ms from now is very
similar to the voltage now. This is why trivial ``just repeat the last value''
prediction is a surprisingly strong baseline, and why you must always report it.
\item \textbf{EEG is dominated by rhythms.} Activity clusters into frequency
bands: delta (0.5--4\,Hz, deep sleep), theta (4--8\,Hz), alpha (8--13\,Hz, relaxed
wakefulness, prominent with eyes closed), beta (13--30\,Hz), gamma ($>$30\,Hz).
Sleep staging is largely reading which rhythms dominate.
\end{itemize}

\subsection{The data we use} \textbf{[established]}

\textbf{Sleep-EDF (sleep-cassette).} A public PhysioNet dataset. \textbf{78
subjects}, most with two full nights recorded. Two EEG channels (Fpz-Cz and
Pz-Oz --- these names are electrode pairs, forehead-to-centre and
parietal-to-occipital), plus an EOG channel that picks up eye movement. Sampled at
\textbf{100\,Hz}, meaning one sample every 10\,milliseconds.

\begin{keybox}
\textbf{Internalise this conversion, you will use it constantly.} At 100\,Hz,
1 sample = 10\,ms. So 127 samples $\approx$ 1.27 seconds. A ``horizon of 16''
means 160\,ms into the future. When we say the model works out to 160\,ms, that
is a horizon of 16 samples.
\end{keybox}

Each night is also \emph{hypnogram-scored}: a human expert labelled every
30-second epoch as wake, one of several non-REM depths, or REM. This is why the
dataset is useful --- it has ground-truth state labels alongside the raw signal.

We also use \textbf{slpdb} (MIT-BIH Polysomnographic Database) for signal
diversity: it carries ECG (heart), blood pressure, respiration and EEG together.
And \textbf{ds005620}, an OpenNeuro propofol-anaesthesia dataset (21 subjects),
used to test whether complexity measures track loss of consciousness.

\section{Part II --- The vocabulary}

\subsection{Forecasting, windows, horizons}

You hand a model a stretch of recent signal --- the \textbf{input window}, length
$W$ --- and ask it to predict the signal at some point in the future. How far
ahead is the \textbf{horizon}, $H$. Both are counted in samples.

\subsection{MSE, and what a number actually means}

Mean squared error. Take the model's prediction, subtract the truth, square it,
average over all cases:
\begin{equation}
\mathrm{MSE} = \frac{1}{n}\sum_{i=1}^{n}\bigl(y_i - \hat{y}_i\bigr)^2
\end{equation}
where $y_i$ is the true value, $\hat{y}_i$ the prediction, $n$ the number of
cases. Lower is better; zero is perfect.

\begin{keybox}
\textbf{Why our MSE numbers hover around 1.0.} We \emph{z-score} every recording
first --- subtract its mean, divide by its standard deviation --- so the signal
has variance 1. That makes numbers comparable across people and machines. The
consequence: \textbf{an MSE of 1.0 means the model is doing no better than
guessing the average.} MSE 0.73 means it explained about 27\,\% of the variance.
So when you see ``GRU 0.7268'', read it as ``about a quarter better than useless''
--- which for brain data at 160\,ms is genuinely respectable.
\end{keybox}

\subsection{The four models we compare}

Listed weakest to strongest. \textbf{You must always report the weak ones} ---
that is the whole discipline.

\begin{itemize}
\item \textbf{Mean-of-trace.} Predict the average of the input window. A floor.
\item \textbf{Persistence.} Predict that the future equals the last value you saw.
Sounds stupid; on autocorrelated signals it is hard to beat. In forecasting
generally this is called the \emph{naive} forecast, and a large literature exists
on methods that fail to beat it.
\item \textbf{Ridge-AR.} A \emph{linear} model. It takes a weighted sum of the
past $W$ samples to predict the future one. ``AR'' is autoregressive --- predicting
a signal from its own past. ``Ridge'' means it is fitted with a penalty on large
weights, which prevents overfitting; the strength of that penalty is called
$\alpha$ and must be tuned. It has a closed-form solution, so it is fast and
exactly optimal for its model class.
\item \textbf{GRU.} Gated Recurrent Unit --- a small neural network that reads the
input sequence one step at a time, keeping an internal memory it chooses what to
write to and forget. It is \emph{nonlinear}, so it can represent relationships a
ridge cannot. This is ``the deep model'' in our comparisons.
\end{itemize}

\begin{keybox}
\textbf{The entire scientific question in our benchmark is:} does the nonlinear
GRU beat the linear ridge, \emph{when the ridge has been tuned properly?} If it
doesn't, the extra complexity bought nothing.
\end{keybox}

\subsection{Train, validation, test --- and why ``subject-held-out''}

You split data three ways. \textbf{Train} fits the model. \textbf{Validation}
picks settings ($\alpha$, network size). \textbf{Test} is touched once, to report
the number.

\textbf{Subject-held-out} means the split is by \emph{person}, not by recording or
window. If any data from a person appears in training, none of that person's data
may appear in test. Otherwise the model can recognise the individual rather than
learn the phenomenon --- EEG is somewhat like a fingerprint. Our split is 39
people train, 12 validation, 27 test, and both nights of a person always travel
together.

\section{Part III --- The statistics, in plain terms}

\subsection{$p$-values}

A $p$-value answers: \emph{if there were genuinely no effect, how often would I see
a result this strong by luck?} $p = 3.7\times10^{-7}$ means roughly once in three
million. It does \textbf{not} say the effect is large or important --- only that it
is unlikely to be noise.

\subsection{Wilcoxon signed-rank test}

We compute each held-out person's error separately, giving 27 paired comparisons
(ridge vs GRU on the same person). The Wilcoxon test asks whether those paired
differences lean consistently one way. We use it instead of a $t$-test because it
makes no assumption that the differences are normally distributed.

\textbf{Say it as:} ``we compared the two models within each person, then asked
whether the direction was consistent across people.''

\subsection{Bootstrap confidence intervals}

Resample your 27 people with replacement, 10{,}000 times, recomputing the average
difference each time. The middle 95\,\% of those values is the confidence
interval. If it excludes zero, the effect is stable. It is a way of asking ``how
much would this move if I'd happened to recruit different people?''

\subsection{Holm correction}

We tested nine horizons. Test enough things and something looks significant by
chance. Holm adjusts the thresholds to control that. Our headline survives it
($p = 2.2\times10^{-6}$ after correction), which is the point of reporting it.

\subsection{Controls}

A control is a deliberately broken version of your experiment that \emph{must}
fail. If it doesn't, your pipeline is manufacturing results. Ours:

\begin{itemize}
\item \textbf{Shuffled targets.} Randomly permute the answers. Both models must
collapse to MSE $\approx$ 1.0. They did (0.991 and 1.022).
\item \textbf{Noise input.} Replace the brain signal with random noise. Same
expectation. They did (0.995 and 0.996).
\item \textbf{AR(2) sanity.} Generate synthetic data where the mathematically
optimal predictor is provably \emph{linear}. The GRU must therefore \emph{fail} to
beat ridge. It did lose ($p = 0.014$). \textbf{This is the strongest control we
have} --- it shows the apparatus cannot manufacture a win for our own model.
\end{itemize}

\section{Part IV --- What ``F9'' means, and every finding}

\textbf{F-numbers are just our internal IDs for findings}, F0 through F9, recorded
in \fp{docs/results/findings-ledger.md}. There is nothing standard about them;
they exist so every claim has an address. \textbf{F9 is the overlap inflation law}
--- the newest result, Part V.

\begin{center}\small
\begin{longtable}{p{0.05\textwidth}p{0.42\textwidth}p{0.45\textwidth}}
\toprule
\textbf{ID} & \textbf{What it is} & \textbf{Verdict, in plain terms} \\
\midrule
F0 & The evaluation apparatus (``gates'' M0--M5) & \textbf{Works.} Passes on planted signals, fails on planted nulls. This is the asset --- a referee we trust. \\
F1 & The 3.116 MSE / $r$=0.972 forecasting result & \textbf{Invalidated.} The overlap illusion. Kept on record as evidence of the bug. \\
F2 & GRU vs trivial baselines on BNCI motor-imagery data & \textbf{Negative} --- but only 2 test subjects and an untuned baseline. Below the standard we now hold ourselves to; needs re-running. \\
F3 & Do heart/breathing signals help predict sleep transitions? & \textbf{Negative (scout).} No detectable extra information. \\
F4 & Do complexity measures beat a simple frequency baseline for wake/NREM/REM? & \textbf{Powered negative.} 78 subjects, 413{,}828 windows. Clearly no. \\
F5 & Same question under propofol anaesthesia & \textbf{Powered negative.} All 21 subjects. Clearly no. \\
F6 & Autonomic signals for micro-arousal prediction & \textbf{Incomplete.} Works on synthetic data; real test needs NSRR data access. \\
F8 & Corrected forecasting benchmark & \textbf{Positive, bounded.} The GRU genuinely beats a tuned ridge out to 160\,ms; the advantage ends by 320\,ms. \\
F9 & The overlap inflation law & \textbf{Exact bound, verified.} 286 configurations, zero violations. \\
\bottomrule
\end{longtable}
\end{center}

\begin{keybox}
\textbf{Seven of nine are negative or incomplete, and that is deliberately on the
record.} The honest summary of this project is: \emph{we built a referee that
bites, pointed it at our own work, and most of what we hoped for did not survive.}
Two things did --- F8 and F9. A reviewer who sees the negatives reported plainly
will trust the positives far more.
\end{keybox}

\section{Part V --- The two real results, worked through}

\subsection{F9: the overlap inflation law} \textbf{[ours]}

\textbf{The mistake.} Build a forecasting example like this: take a stretch of
signal, use the first $W$ samples as input $X$, and ask the model to predict a
window $Y$ that starts $H$ samples later:
\begin{equation}
X = \mathrm{seq}[0 : W], \qquad Y = \mathrm{seq}[H : H+W].
\end{equation}
Both are length $W$. Now look at what is actually in $Y$. Its $k$-th entry is
$Y[k] = \mathrm{seq}[H+k]$. As long as $H+k < W$, that value is \emph{also inside
$X$}, at position $H+k$. So:

\begin{itemize}
\item $W - H$ entries of $Y$ are \textbf{copyable} --- already handed to the model.
\item Only $H$ entries are genuinely in the future.
\end{itemize}

With $W = 127$ and $H = 1$: 126 copyable, 1 real. The model learned to shift its
input by one position and got 126 of 127 ``predictions'' exactly right. Scored
over the whole window, it looked near-perfect.

\textbf{The derivation.} Score \emph{one} model's \emph{one} set of predictions
three ways, changing only which positions you average over:
\begin{align}
\mathrm{MSE}_{\text{overlap}} &= \text{average over all } W \text{ positions (what the flawed protocol reports)}\\
\mathrm{MSE}_{\text{copy}} &= \text{average over the } W-H \text{ copyable positions}\\
\mathrm{MSE}_{\text{future}} &= \text{average over the } H \text{ genuinely-future positions (the honest number)}
\end{align}
Because an average over a set equals the weighted average of its parts, these are
related \emph{by definition}, not by any assumption:
\begin{equation}
\mathrm{MSE}_{\text{overlap}}
= \frac{(W-H)\,\mathrm{MSE}_{\text{copy}} + H\,\mathrm{MSE}_{\text{future}}}{W}
\label{eq:identity}
\end{equation}
Define the inflation factor as how much better the flawed score looks than the
honest one, and substitute \eqref{eq:identity}:
\begin{equation}
\boxed{\;
\text{inflation} \;=\; \frac{\mathrm{MSE}_{\text{future}}}{\mathrm{MSE}_{\text{overlap}}}
\;=\;
\frac{W/H}{\,1 + \frac{W-H}{H}\cdot\frac{\mathrm{MSE}_{\text{copy}}}{\mathrm{MSE}_{\text{future}}}\,}
\;\le\; \frac{W}{H}\;}
\label{eq:law}
\end{equation}

\textbf{Reading equation \eqref{eq:law} in words.} The denominator is at least 1,
so the whole thing is at most $W/H$. If the model copies the copyable positions
perfectly, $\mathrm{MSE}_{\text{copy}} = 0$, the denominator becomes exactly 1, and
the inflation is exactly $W/H$. If it copies them badly, the denominator grows and
the inflation is smaller. \textbf{So $W/H$ is a ceiling, and a model that is good
at copying sits right on it.}

\textbf{Why this matters to other people.} $W$ and $H$ are printed in the methods
section of any paper. You do not need their data, their code, or their model.
Divide the two numbers and you have an upper bound on how inflated their result
is.

\textbf{Verification.} 286 combinations of $W$ and $H$, across 11 signals: three
Sleep-EDF channels, four slpdb channels, and four synthetic processes. Result:
\textbf{zero violations}, median observed/predicted $= 1.000$. Our own historical
bug ($W=127$, $H=1$) predicted $127\times$; we had measured $126.1\times$.

\begin{keybox}
\textbf{The single best piece of evidence:} we ran it on \emph{pure random noise},
where nothing whatsoever is predictable. It still inflated by exactly the
predicted factor. That is what proves this is a property of the \emph{scoring
rule}, not of brain data --- the copyable positions are sitting in the input
regardless of whether anything is forecastable.
\end{keybox}

\subsection{F8: the corrected benchmark} \textbf{[ours]}

Having fixed the scoring, we re-ran the actual experiment: 78 subjects, 39/12/27
split, predicting a single strictly-future sample so that input and target share
nothing.

\begin{center}\small
\begin{tabular}{rrrrrr}
\toprule
Horizon & Time & Persistence & Ridge-AR & GRU & Holm $p$ \\
\midrule
4 & 40\,ms & 0.5691 & 0.4399 & \textbf{0.3950} & $1.3\times10^{-7}$ \\
16 & 160\,ms & 1.3814 & 0.8298 & \textbf{0.7268} & $2.2\times10^{-6}$ \\
32 & 320\,ms & 1.7600 & \textbf{0.8700} & 0.8756 & n.s. \\
256 & 2560\,ms & 1.8746 & \textbf{0.9942} & 1.0099 & $1.3\times10^{-7}$ \\
\bottomrule
\end{tabular}
\end{center}

At 160\,ms the GRU beats a fully-tuned ridge on 26 of 27 held-out people, and is
47\,\% better than persistence. By 320\,ms the advantage is gone and never returns.

\textbf{A second lesson hidden in this table.} Early on, we ran this with the
ridge's $\alpha$ fixed at an arbitrary value, and later with a search grid that
was too narrow. \emph{Both gave the wrong scientific conclusion, in opposite
directions.} Only after extending the grid until the search stopped picking its
own edge did the answer stabilise. \textbf{An under-tuned baseline is the second
way these results inflate}, and we caught ourselves doing it twice.

\textbf{What broke under scrutiny.} We had reported that ridge beats the GRU at
640\,ms ($p=4.1\times10^{-5}$). Under a stricter windowing check that became a tie
($p = 0.58$). We narrowed the claim rather than removing the test. Say this out
loud when asked --- it is the most credible thing in the project.

\section{Part VI --- Every paper, and what each is for}

\subsection{Establishing that our overlap axis is new}

These all address leakage --- and all address a \emph{different} kind than ours.
They are how we know the finding is novel.

\begin{itemize}
\item \textbf{Bernett et al.\ (2024), \emph{Nature Methods}.} Seven questions to
prevent data leakage in biological machine learning. \emph{Why it matters:} the
authoritative checklist --- and none of its seven questions asks whether a
regression target overlaps its own input.
\item \textbf{Del Pup et al.\ (2025), \emph{Computers in Biology and Medicine}.}
Compared five cross-validation schemes across over 100{,}000 trained EEG models.
\emph{Why it matters:} the most thorough study of EEG splitting that exists, and
our artifact survives every scheme in it.
\item \textbf{Apicella et al.\ (2024), \emph{AI Review}.} A taxonomy of leakage
across the ML workflow. \emph{Why it matters:} the within-example case is absent.
\item \textbf{Lee et al.\ (2023), \emph{Scientific Reports}.} Overlapping
cross-validation inflating EEG diagnosis performance. \emph{Why:} overlap
\emph{between} augmented samples --- the other axis.
\item \textbf{Dehghani et al.\ (2019), \emph{Sensors}.} Overlapping vs
non-overlapping sliding windows in activity recognition; the apparent benefit of
overlap was an artifact of subject-dependent validation. \emph{Why:} same lesson,
still the between-example axis.
\item \textbf{Wosiak et al.\ (2025), \emph{Algorithms}.} Argues window overlap
should be a first-class reported parameter --- and still never mentions horizon
versus window length.
\end{itemize}

\subsection{Establishing that the paper's shape is respected}

\begin{itemize}
\item \textbf{McGreivy \& Hakim (2024), \emph{Nature Machine Intelligence}.}
Reviewed machine-learning papers claiming to beat numerical solvers and found
\textbf{79\,\% (60/76) compared against a weak baseline}. \emph{Why it matters:}
this is the template for our paper, in a different field, in a top journal. It
proves the genre is publishable at the highest level.
\item \textbf{Lin (2019), \emph{SIGIR Forum}} and \textbf{Yang et al.\ (2019).}
The ``neural hype'' papers in information retrieval; re-ran neural models against
strong baselines and found most reported gains illusory. \emph{Why:} the same
move, a decade earlier, in another field.
\item \textbf{Beck et al.\ (2025), \emph{Applied Intelligence}.} ``Mind the naive
forecast!'' --- on volatile series no method reliably beat the no-change forecast.
\emph{Why:} justifies always reporting persistence.
\item \textbf{Wu et al.\ (2025), \emph{IEEE Access}.} Reviewed 149 EEG
seizure-prediction studies and named \textbf{``few linear baselines''} as a
field-wide failure, explicitly calling for benchmarks against well-tuned linear
models. \emph{Why it matters:} our corrected benchmark answers a gap the field has
already stated in print. Quote this one.
\item \textbf{Puffay et al.\ (2023), \emph{J.\ Neural Engineering}.} 29 EEG-speech
deep learning studies; recurring pitfalls were biased cross-validation and
leakage.
\end{itemize}

\subsection{Direct comparators --- know these or you will be caught out}

\begin{itemize}
\item \textbf{Pankka et al.\ (2024).} WaveNet forecasting resting-state EEG
\textbf{150\,ms ahead}, beating an autoregressive model on amplitude and phase.
\emph{Why it matters:} \textbf{this is our exact horizon regime.} ``A deep model
beats AR at $\sim$150\,ms'' is already published. Our contribution therefore
\emph{cannot} be that claim --- it is the audit, the law, and the horizon-resolved
crossover. Never claim novelty for beating AR at 150\,ms.
\item \textbf{Onojima et al.\ (2025).} Predicts the complex analytic signal from
raw \emph{multichannel} EEG, beating AR both within and across subjects.
\emph{Why:} a comparator, and evidence that multichannel is where a nonlinear
model should earn its keep --- our single-channel setup is the likeliest reason our
advantage closes so early.
\end{itemize}

\subsection{Background for the negative findings}

\begin{itemize}
\item \textbf{Casali et al.\ (2013)} --- the Perturbational Complexity Index (PCI),
which perturbs cortex with TMS and measures the complexity of the response as a
consciousness marker. \emph{Why it matters:} F4 and F5 test whether \emph{passive}
complexity (no perturbation) can do something similar. They cannot.
\textbf{Careful:} we have never measured PCI and must never claim to have beaten
it.
\item \textbf{Lempel--Ziv complexity, permutation entropy} (Bandt \& Pompe),
\textbf{multiscale entropy} --- the specific complexity measures in F4/F5. All are
ways of asking ``how compressible / how disordered is this signal?''
\item \textbf{McAllester \& Stratos (2018).} Proves you cannot obtain a
distribution-free high-confidence lower bound on mutual information much larger
than $\log N$ from $N$ samples. \emph{Why it matters:} it is why our RFS score
(below) is described as a bounded functional of held-out predictions and
\textbf{not} as a mutual-information estimate. Getting this wrong in a talk would
be a real error.
\end{itemize}

\subsection{One more equation you will be asked about: RFS}

Residual Forecastability Score, used in the F3--F6 gates:
\begin{equation}
\mathrm{RFS}_{\text{bits}} = \frac{\mathrm{NLL}_{\text{baseline}} - \mathrm{NLL}_{\text{model}}}{\ln 2}
\end{equation}
NLL is negative log-likelihood --- roughly, how surprised a model is by the truth
(lower is better). So the numerator is how much less surprised the full model is
than the baseline. Dividing by $\ln 2$ converts to bits. \textbf{Positive means
the extra information helped; zero or negative means it did not.} Our real-data
values came out negative, which is why F3--F5 are negatives.

\section{Part VII --- Glossary}

\begin{center}\small
\begin{longtable}{p{0.22\textwidth}p{0.71\textwidth}}
\toprule
\textbf{Term} & \textbf{Meaning} \\
\midrule
Autocorrelation & How similar a signal is to itself a moment later. High in EEG. \\
Horizon ($H$) & How far ahead you predict, in samples. $H=16$ at 100\,Hz is 160\,ms. \\
Window ($W$) & How many samples you feed in, or score over. \\
Epoch & (a) A 30-second scored segment of sleep; (b) one pass through training data. Context decides. \\
Hypnogram & The expert's sleep-stage labelling of a night. \\
z-scoring & Rescaling to mean 0, variance 1, so MSE 1.0 means ``no better than the mean''. \\
Ridge / $\alpha$ & Linear regression with a penalty on large weights; $\alpha$ is the penalty strength, tuned on validation. \\
GRU & Gated Recurrent Unit; small nonlinear sequence network with learned memory. \\
Persistence & Predict that nothing changes. The naive baseline. \\
Subject-held-out & Splitting by person, so no individual appears in both train and test. \\
Leakage & Information reaching the model that it should not have; inflates results. \\
Bootstrap & Resampling your data many times to see how stable a number is. \\
Wilcoxon & A paired test that does not assume normally distributed differences. \\
Holm & A correction for having run many tests. \\
Control & A deliberately broken run that must fail; proves the pipeline is honest. \\
RFS & Residual Forecastability Score, in bits. Positive means extra signal helped. \\
PCI & Perturbational Complexity Index --- a TMS-based consciousness measure. \textbf{We do not measure it.} \\
Gate & Our pass/fail rule for whether a result may be claimed. \\
\bottomrule
\end{longtable}
\end{center}

\section{Part VIII --- Questions you will be asked}

\textbf{``What's the contribution?''} A closed-form bound letting anyone compute
how inflated a published forecasting number is from two numbers in its methods
section, plus a corrected benchmark showing where a neural model genuinely beats a
classical one.

\textbf{``Isn't this just data leakage?''} Different axis. Standard leakage is
information crossing the train/test split, fixed by splitting correctly. Ours is
inside a single example --- the answer overlaps the question --- so splitting
cannot fix it.

\textbf{``How is the law not just a curve fit?''} It's derived from the definition
of an average, then checked 286 times without a violation. On pure noise, where
nothing is predictable, it still holds exactly.

\textbf{``Does your model actually work?''} Out to 160\,ms, yes, on 26 of 27
held-out people against a properly tuned baseline. Past 320\,ms, no. We also had a
long-horizon claim that failed our own check and we reported that.

\textbf{``Is this clinically useful?''} Not yet, and we don't claim it. The
timescale that matters for predicting state changes is around a minute --- hundreds
of times longer than what we demonstrate.

\textbf{``What's next?''} A prevalence survey: how often does published work use a
horizon smaller than its window? That converts a tool into a finding about the
field, and it is the McGreivy move.

\textbf{``What are you least sure about?''} Single channel, single corpus. The
advantage may close early precisely because one electrode has no spatial
information to exploit --- Onojima's multichannel result suggests that is where to
look next.

\vspace{6pt}\hrule\vspace{5pt}
\noindent\footnotesize
\textbf{Source documents.} \fp{docs/results/findings-ledger.md} (all findings)
$\cdot$ \fp{docs/results/overlap_inflation_law_2026-07-28.md} (F9) $\cdot$
\fp{docs/results/forecast_crossover_2026-07-27.md} (F8) $\cdot$
\fp{docs/research/baseline_fairness_literature_2026-07-27.md} (citations)
$\cdot$ \fp{docs/research/eval_leakage_audit_2026-07-21.md} (the original
bug audit).

\end{document}
